\documentclass{article}
\usepackage[utf8]{inputenc}
\usepackage[vietnamese,english]{babel}
\usepackage{times}
\usepackage{amsmath,amssymb}
\usepackage{tabularx}
\usepackage{array}
\usepackage{booktabs}
\usepackage{graphicx}
\usepackage{xcolor}
\usepackage[
  paperwidth=29.7cm,
  paperheight=16.706cm,
  top=1.0cm,
  bottom=1.0cm,
  left=1.2cm,
  right=1.2cm,
  footskip=0pt,
  headheight=0pt,
  headsep=0pt
]{geometry}

\pagestyle{empty}
\setlength{\parindent}{0pt}
\setlength{\parskip}{0pt}

% Định nghĩa màu sắc thương hiệu
\definecolor{brandGreen}{RGB}{46, 125, 50}
\definecolor{brandRed}{RGB}{211, 47, 47}
\definecolor{brandBlue}{RGB}{21, 101, 192}

\begin{document}
\selectlanguage{vietnamese}

% ==============================================================================
% SLIDE 1: TRANG TIÊU ĐỀ
% ==============================================================================
\begin{center}
\vspace*{0.6cm}
{\fontsize{24}{30}\selectfont \textbf{TỔ KHOA HỌC TỰ NHIÊN}}\\[0.6cm]
{\fontsize{34}{42}\selectfont \textbf{\color{brandGreen}HƯỚNG DẪN ÔN TẬP GIỮA HỌC KỲ I}}\\[0.4cm]
{\fontsize{28}{36}\selectfont \textbf{MÔN: KHOA HỌC TỰ NHIÊN 9 (HÓA HỌC \& VẬT LÍ)}}\\[0.4cm]
{\fontsize{22}{28}\selectfont \textbf{NĂM HỌC: 2026 - 2027}}\\[0.8cm]
{\fontsize{24}{32}\selectfont \textit{Phân môn: Hóa học (Hợp chất hữu cơ \& Alkane) -- Vật lí (Năng lượng \& Khúc xạ ánh sáng)}}\\[1.0cm]
{\fontsize{24}{30}\selectfont \textbf{\color{brandRed}Thầy TRẦN VĂN THẠNH -- 0777.470.803}}\\[0.3cm]
{\fontsize{20}{26}\selectfont CS1: 6/15 Nguyễn Hoàng, P. Kim Long, TP Huế \quad | \quad CS2: 24 Đặng Thái Thân, TP Huế}
\end{center}
\newpage

% ==============================================================================
% PHẦN I: HÓA HỌC - TỰ LUẬN
% ==============================================================================

% SLIDE 2: HÓA TỰ LUẬN CÂU 1
\fontsize{24}{32}\selectfont
\textbf{Câu 1 (Hóa học).} Hãy sắp xếp các chất sau vào các cột thích hợp trong bảng:\\
\centerline{$\text{C}_6\text{H}_6,\ \text{CaCO}_3,\ \text{C}_4\text{H}_{10},\ \text{C}_2\text{H}_6\text{O},\ \text{NaNO}_3,\ \text{CH}_3\text{NO}_2,\ \text{NaHCO}_3,\ \text{C}_2\text{H}_3\text{O}_2\text{Na}$}
\vspace{12pt}

\begin{center}
\begin{tabular}{|>{\centering\arraybackslash}p{8.5cm}|>{\centering\arraybackslash}p{9.0cm}|>{\centering\arraybackslash}p{8.5cm}|}
\hline
\multicolumn{2}{|c|}{\textbf{HỢP CHẤT HỮU CƠ}} & \textbf{HỢP CHẤT VÔ CƠ}\\[4pt]
\hline
\textbf{Hydrocarbon} & \textbf{Dẫn xuất của Hydrocarbon} & \textbf{Hợp chất vô cơ}\\[4pt]
\hline
\vspace{35pt} & \vspace{35pt} & \vspace{35pt}\\
\hline
\end{tabular}
\end{center}
\vspace*{\fill}
\newpage

% SLIDE 3: HÓA TỰ LUẬN CÂU 2
\fontsize{25}{33}\selectfont
\textbf{Câu 2 (Hóa học).} Hãy viết công thức cấu tạo đầy đủ và thu gọn có thể có của các chất có công thức phân tử sau:
\vspace{16pt}

\begin{tabularx}{\linewidth}{@{}X X@{}}
a) $\text{C}_2\text{H}_5\text{Br}$ & b) $\text{C}_4\text{H}_{10}$\\[16pt]
c) $\text{C}_3\text{H}_6$ & d) $\text{C}_3\text{H}_8\text{O}$
\end{tabularx}
\vspace*{\fill}
\newpage

% SLIDE 4: HÓA TỰ LUẬN CÂU 3
\fontsize{26}{34}\selectfont
\textbf{Câu 3 (Hóa học).} Gas dùng đun nấu có chứa alkane chủ yếu nào? Viết phương trình hoá học biểu diễn phản ứng cháy của chúng.
\vspace*{\fill}
\newpage

% SLIDE 5: HÓA TỰ LUẬN CÂU 4
\fontsize{26}{34}\selectfont
\textbf{Câu 4 (Hóa học).} Alkane là gì? Viết công thức chung của alkane. Nêu các ứng dụng quan trọng của alkane trong đời sống và sản xuất.
\vspace*{\fill}
\newpage

% SLIDE 6: HÓA TỰ LUẬN CÂU 5
\fontsize{25}{33}\selectfont
\textbf{Câu 5 (Hóa học).} Alkane X có khối lượng phân tử $44\text{ amu}$.\\
a. Xác định công thức phân tử của X và viết công thức cấu tạo có thể có ứng với CTPT vừa tìm được.\\
b. Đốt cháy hoàn toàn X trong không khí rồi dẫn sản phẩm cháy qua nước vôi trong dư. Nêu hiện tượng và viết PTHH xảy ra.
\vspace*{\fill}
\newpage

% ==============================================================================
% PHẦN I: HÓA HỌC - TRẮC NGHIỆM
% ==============================================================================

% SLIDE 7: HÓA TRẮC NGHIỆM CÂU 1 & CÂU 2
\fontsize{24}{31}\selectfont
\textbf{Câu 1.} Chất nào sau đây không thuộc loại chất hữu cơ?
\vspace{4pt}

\begin{tabularx}{\linewidth}{@{}X X X X@{}}
\textbf{A.} $\text{CH}_4$. & \textbf{B.} $\text{CH}_3\text{Cl}$. & \textbf{C.} $\text{CH}_3\text{COONa}$. & \textbf{D.} $\text{CO}_2$.
\end{tabularx}

\vspace{10pt}

\textbf{Câu 2.} Cho biết công thức cấu tạo dưới đây là của hợp chất có CTPT nào?
\vspace{2pt}

\begin{center}
\includegraphics[width=6.0cm]{c:/Antigravity_Thanh/San_Pham/images_khtn9/khtn9_butane_structure.png}
\end{center}
\vspace{2pt}

\begin{tabularx}{\linewidth}{@{}X X X X@{}}
\textbf{A.} $\text{C}_4\text{H}_8$. & \textbf{B.} $\text{C}_5\text{H}_{10}$. & \textbf{C.} $\text{C}_4\text{H}_{10}$. & \textbf{D.} $\text{C}_6\text{H}_{10}$.
\end{tabularx}
\vspace*{\fill}
\newpage

% SLIDE 8: HÓA TRẮC NGHIỆM CÂU 4
\fontsize{25}{33}\selectfont
\textbf{Câu 4.} Cho nhiệt lượng tỏa ra khi đốt cháy hoàn toàn $1\text{ mol}$ các chất ethane, propane, butane và pentane ($\text{C}_5\text{H}_{12}$) lần lượt là $1570\text{ kJ/mol}$; $2220\text{ kJ/mol}$; $2875\text{ kJ/mol}$ và $3536\text{ kJ/mol}$. Khi đốt cháy hoàn toàn $1\text{ gam}$ chất nào sẽ thu được nhiệt lượng lớn nhất?
\vspace{16pt}

\begin{tabularx}{\linewidth}{@{}X X X X@{}}
\textbf{A.} Ethane. & \textbf{B.} Propane. & \textbf{C.} Pentane. & \textbf{D.} Butane.
\end{tabularx}
\vspace*{\fill}
\newpage

% SLIDE 9: HÓA TRẮC NGHIỆM CÂU 5
\fontsize{24}{32}\selectfont
\textbf{Câu 5.} Trong số các câu phát biểu sau đây, số câu đúng là:\\
(1) Mỗi CTPT có thể biểu diễn cho nhiều chất hữu cơ.\\
(2) Hydrocarbon là hợp chất mà thành phần phân tử chỉ chứa các nguyên tố carbon, hydrogen.\\
(3) Mỗi hợp chất hữu cơ có thể có nhiều CTPT.\\
(4) Mỗi hợp chất hữu cơ có thể có nhiều CTCT.
\vspace{14pt}

\begin{tabularx}{\linewidth}{@{}X X X X@{}}
\textbf{A.} $2$. & \textbf{B.} $3$. & \textbf{C.} $4$. & \textbf{D.} $1$.
\end{tabularx}
\vspace*{\fill}
\newpage

% SLIDE 10: HÓA ĐÚNG/SAI
\fontsize{24}{32}\selectfont
\textbf{Phần Trắc nghiệm Đúng -- Sai (Hóa học).} Công thức phân tử cho biết thành phần nguyên tố và số lượng nguyên tử của mỗi nguyên tố trong hợp chất hữu cơ:
\vspace{12pt}

a) Công thức phân tử của ethane là $\text{CH}_4$. \hfill $[\quad]$\\[10pt]
b) Công thức phân tử cho biết cách các nguyên tử liên kết với nhau trong phân tử. \hfill $[\quad]$\\[10pt]
c) Công thức cấu tạo biểu diễn liên kết giữa các nguyên tử trong phân tử. \hfill $[\quad]$\\[10pt]
d) Để hiểu rõ cấu trúc của một hợp chất hữu cơ, cần kết hợp cả công thức phân tử và công thức cấu tạo. \hfill $[\quad]$
\vspace*{\fill}
\newpage

% ==============================================================================
% PHẦN II: VẬT LÍ - TRẮC NGHIỆM & BÀI TẬP
% ==============================================================================

% SLIDE 11: VẬT LÍ CÂU 1 & CÂU 2
\fontsize{25}{33}\selectfont
\textbf{Câu 1.} Trong các nguồn năng lượng dưới đây, nguồn năng lượng nào là năng lượng tái tạo?
\vspace{6pt}

\begin{tabularx}{\linewidth}{@{}X X X X@{}}
\textbf{A.} Than đá. & \textbf{B.} Khí tự nhiên. & \textbf{C.} Dầu mỏ. & \textbf{D.} Năng lượng gió.
\end{tabularx}

\vspace*{\fill}

\textbf{Câu 2.} Khi một tia sáng truyền từ không khí vào nước, phát biểu nào sau đây đúng theo định luật khúc xạ ánh sáng?
\vspace{6pt}

\textbf{A.} Tia khúc xạ luôn vuông góc với mặt phân cách.\\
\textbf{B.} Tia tới, tia khúc xạ và pháp tuyến tại điểm tới cùng nằm trong một mặt phẳng.\\
\textbf{C.} Góc khúc xạ luôn bằng góc tới.\\
\textbf{D.} Tia sáng luôn bị phản xạ hoàn toàn.
\vspace*{\fill}
\newpage

% SLIDE 12: VẬT LÍ CÂU 3 & CÂU 4
\fontsize{25}{33}\selectfont
\textbf{Câu 3.} Một máy nâng hàng có công suất $400\text{ W}$, thực hiện công $72\,000\text{ J}$ để nâng các kiện hàng lên cao. Thời gian máy thực hiện công là
\vspace{6pt}

\begin{tabularx}{\linewidth}{@{}X X X X@{}}
\textbf{A.} $120\text{ s}$. & \textbf{B.} $160\text{ s}$. & \textbf{C.} $180\text{ s}$. & \textbf{D.} $240\text{ s}$.
\end{tabularx}

\vspace*{\fill}

\textbf{Câu 4.} Trong buổi lao động, một học sinh tác dụng lực kéo $80\text{ N}$ theo phương chuyển động để kéo xe chở dụng cụ đi được quãng đường $15\text{ m}$. Công của lực kéo là
\vspace{6pt}

\begin{tabularx}{\linewidth}{@{}X X X X@{}}
\textbf{A.} $1\,200\text{ J}$. & \textbf{B.} $1\,200\text{ W}$. & \textbf{C.} $95\text{ J}$. & \textbf{D.} $5,3\text{ J}$.
\end{tabularx}
\vspace*{\fill}
\newpage

% SLIDE 13: VẬT LÍ ĐÚNG/SAI CÂU 5
\fontsize{24}{32}\selectfont
\textbf{Câu 5.} Nhóm học sinh khảo sát chuyển động của quả bóng khối lượng $0,4\text{ kg}$ thả rơi từ độ cao $5\text{ m}$. Bỏ qua sức cản không khí, gốc thế năng ở mặt đất:
\vspace{12pt}

a) Thế năng trọng trường của bóng phụ thuộc vào khối lượng và độ cao của bóng so với mốc thế năng. \hfill $[\quad]$\\[10pt]
b) Cơ năng của quả bóng bằng tổng động năng và thế năng trọng trường của nó. \hfill $[\quad]$\\[10pt]
c) Khi quả bóng ở độ cao $2\text{ m}$ so với mặt đất, thế năng trọng trường của bóng là $8\text{ J}$. \hfill $[\quad]$\\[10pt]
d) Khi quả bóng ở độ cao $2\text{ m}$, động năng của bóng bằng $8\text{ J}$. \hfill $[\quad]$
\vspace*{\fill}
\newpage

% SLIDE 14: VẬT LÍ ĐÚNG/SAI CÂU 6
\fontsize{24}{32}\selectfont
\textbf{Câu 6.} Chiếu một tia sáng từ không khí vào khối thủy tinh có chiết suất $1,5$ (chiết suất không khí xấp xỉ 1), góc tới $i = 30^\circ$:
\vspace{12pt}

a) Góc tới là góc hợp bởi tia tới và pháp tuyến tại điểm tới. \hfill $[\quad]$\\[10pt]
b) Chiết suất tuyệt đối của thủy tinh là chiết suất tỉ đối của thủy tinh đối với chân không. \hfill $[\quad]$\\[10pt]
c) Giá trị của $\sin r$ khi tia sáng truyền vào thủy tinh bằng khoảng $0,33$. \hfill $[\quad]$\\[10pt]
d) Nếu tăng góc tới lên $45^\circ$ thì góc khúc xạ trong thủy tinh cũng bằng $45^\circ$. \hfill $[\quad]$
\vspace*{\fill}
\newpage

% SLIDE 15: VẬT LÍ TRẢ LỜI NGẮN CÂU 7
\fontsize{24}{32}\selectfont
\textbf{Câu 7.} Một vật khối lượng $2\text{ kg}$ thả nhẹ từ vị trí A cao $5\text{ m}$ so với mặt đất. Xét vị trí B cao $3\text{ m}$. Bỏ qua sức cản không khí, chọn mốc thế năng ở mặt đất:
\vspace{10pt}

a) Đơn vị của thế năng trong hệ SI là gì? \hfill $\dots\dots\dots\dots\dots\dots$\\[8pt]
b) Khi vật rơi xuống, thế năng trọng trường chuyển hóa thành dạng năng lượng nào? \hfill $\dots\dots\dots\dots\dots\dots$\\[8pt]
c) Thế năng trọng trường của vật tại vị trí B bằng bao nhiêu jun? \hfill $\dots\dots\dots\dots\dots\dots$\\[8pt]
d) Động năng của vật tại vị trí B bằng bao nhiêu jun? \hfill $\dots\dots\dots\dots\dots\dots$
\vspace*{\fill}
\newpage

% SLIDE 16: VẬT LÍ TỰ LUẬN CÂU 8
\fontsize{24}{32}\selectfont
\textbf{Câu 8.} Xe điện vận chuyển hàng hóa chịu lực kéo $600\text{ N}$ cùng hướng chuyển động. Ban đầu chuyển động với vận tốc $3\text{ m/s}$ trong $4\text{ phút}$.\\
a. Tính quãng đường xe đi được và công của lực kéo thực hiện.\\
b. Người điều khiển tăng vận tốc lên $5\text{ m/s}$ trong khi vẫn giữ nguyên lực kéo. Trong cùng khoảng thời gian $4\text{ phút}$, tính công của lực kéo và công suất của xe. So sánh công suất trước và sau khi tăng vận tốc.
\vspace*{\fill}
\newpage

% SLIDE 17: VẬT LÍ TỰ LUẬN CÂU 9
\fontsize{25}{33}\selectfont
\textbf{Câu 9.} Chiếu tia sáng từ không khí vào một khối thủy tinh trong suốt có chiết suất $1,5$ (chiết suất không khí xấp xỉ 1). Góc tới ban đầu bằng $45^\circ$.\\
a. Áp dụng định luật khúc xạ ánh sáng, tính góc khúc xạ trong thủy tinh.\\
b. Tính tốc độ ánh sáng đi trong môi trường thủy tinh (biết $c = 3.10^8\text{ m/s}$).
\vspace*{\fill}
\newpage

% ==============================================================================
% SLIDE 18: BẢNG TỔNG HỢP ĐÁP ÁN TOÀN BỘ ĐỀ CƯƠNG
% ==============================================================================
\begin{center}
{\fontsize{24}{30}\selectfont \textbf{\color{brandGreen}BẢNG TỔNG HỢP ĐÁP ÁN ĐỀ CƯƠNG ÔN TẬP KHTN 9}}\\[0.4cm]
\end{center}

\fontsize{18}{24}\selectfont
\textbf{1. Trắc nghiệm nhiều lựa chọn:}\\
\begin{tabularx}{\linewidth}{|c|*{4}{>{\centering\arraybackslash}X|}|c|*{4}{>{\centering\arraybackslash}X|}}
\hline
\textbf{Hóa} & \textbf{1} & \textbf{2} & \textbf{4} & \textbf{5} & \textbf{Lí} & \textbf{1} & \textbf{2} & \textbf{3} & \textbf{4}\\
\hline
\textbf{Đ/A} & \textbf{\color{brandRed}D} & \textbf{\color{brandRed}C} & \textbf{\color{brandRed}A} & \textbf{\color{brandRed}A} & \textbf{Đ/A} & \textbf{\color{brandRed}D} & \textbf{\color{brandRed}B} & \textbf{\color{brandRed}C} & \textbf{\color{brandRed}A}\\
\hline
\end{tabularx}
\vspace{8pt}

\textbf{2. Trắc nghiệm Đúng / Sai:}\\
$\bullet$ \textbf{Hóa học:} a - \textbf{Sai}, b - \textbf{Sai}, c - \textbf{Đúng}, d - \textbf{Đúng}.\\
$\bullet$ \textbf{Vật lí (Câu 5):} a - \textbf{Đúng}, b - \textbf{Đúng}, c - \textbf{Đúng}, d - \textbf{Sai}.\\
$\bullet$ \textbf{Vật lí (Câu 6):} a - \textbf{Đúng}, b - \textbf{Đúng}, c - \textbf{Đúng}, d - \textbf{Sai}.
\vspace{6pt}

\textbf{3. Trắc nghiệm Trả lời ngắn \& Bài toán Vật lí:}\\
$\bullet$ \textbf{Câu 7:} a) \textbf{Jun (J)} \quad | \quad b) \textbf{Động năng} \quad | \quad c) \textbf{60 J} \quad | \quad d) \textbf{40 J}.\\
$\bullet$ \textbf{Câu 8:} a) $s = \mathbf{720\text{ m}},\ A = \mathbf{432\text{ kJ}}$ \quad | \quad b) $s = \mathbf{1200\text{ m}},\ A = \mathbf{720\text{ kJ}},\ \mathcal{P}_2 = \mathbf{3000\text{ W}}$ (tăng $1200\text{ W}$, gấp $1,67$ lần).\\
$\bullet$ \textbf{Câu 9:} a) $r \approx \mathbf{28,13^\circ}$ \quad | \quad b) $v = \mathbf{2.10^8\text{ m/s}}$.

\end{document}
